% !PRINTED 475
\weberpage{475}

%<<CONT
which is again the identical permutation. \gesperrt{Thus the permutation
opposite to $\pi_a^{-1}$ is again the original permutation $\pi_a$.}

If $\pi_a \pi_b = \pi_c$, then
\[
\pi_c^{-1} = \pi_b^{-1} \pi_a^{-1},
\]
for $\pi_a \pi_b \pi_b^{-1} \pi_a^{-1}$ is the identical
permutation.

If we denote the identical permutation by $\pi_0$, then

\weq{10}{\pi_a \pi_a^{-1} = \pi_a^{-1} \pi_a = \pi_0,}

and $\pi_0$ is opposite to itself (which does not exclude the
possibility that there are still other permutations opposite to
themselves). By composition with the identical permutation no other
permutation is changed; thus

\weq{11}{\pi_0 \pi_a = \pi_a \pi_0 = \pi_a.}

\gesperrt{2. If $\pi_c$ is composed of $\pi_a$ and $\pi_b$, that is,}

\weq{12}{\pi_a \pi_b = \pi_c,}

\gesperrt{then not only is $\pi_c$ completely and uniquely determined
by $\pi_a$ and $\pi_b$, but also $\pi_a$ by $\pi_b$ and
$\pi_c$, and $\pi_b$ by $\pi_a$ and $\pi_c$.}

For it follows from (12) that

\weq{13}{\pi_a = \pi_c \pi_b^{-1}, \quad \pi_b = \pi_a^{-1} \pi_c.}

\gesperrt{3. A system
\[
Q = \pi, \pi', \pi'', \pi''' \ldots,
\]
selected from the totality of all $\Pi(m)$ permutations of $m$
elements, which is so closed and complete in itself that any two, and
consequently also arbitrarily many, permutations of $Q$ always yield
by composition a permutation of the same system, is called a
permutation group. By the degree of a permutation group we understand
the number of permutations that it contains.}

The totality of all $\Pi(m)$ permutations forms such a group.
Likewise the single identical permutation is a group by itself. What
may still lie between these two extreme cases of groups is what will
occupy us in the following sections.
